\documentclass[10pt,journal,compsoc]{IEEEtran}
\usepackage{amsmath,amssymb,booktabs,multirow,graphicx,subcaption,url,xcolor,threeparttable}
\usepackage[hidelinks]{hyperref}
\graphicspath{{paper_assets/}}
\newcommand{\AUC}{\operatorname{AUC}}
\newcommand{\EER}{\operatorname{EER}}
\newcommand{\TAR}{\operatorname{TAR}}
\newcommand{\FAR}{\operatorname{FAR}}
\title{Parameter-Efficient Periocular Fusion for Cross-Condition and Cross-Dataset Occlusion-Robust Face Verification}
\author{Pararajasingham Sangiya\thanks{Independent Researcher. ORCID: 0009-0005-3407-2154.}}
\begin{document}
\maketitle
\begin{abstract}
Face verification under masks and sunglasses requires identity-discriminative representations when part of the face is unavailable or visually altered. This paper studies whether parameter-efficient adaptation and explicit periocular information improve verification without target-domain fine-tuning. We use DINOv2 representations with low-rank adaptation (LoRA), a whole-face branch, a periocular branch, cosine similarity, and fixed or learned fusion. The evaluation is identity-disjoint and verification-centered: identity classification is used only on training identities, while test identities are evaluated through genuine and impostor similarity scores. We compare five configurations: frozen whole-face DINOv2, whole-face LoRA, periocular LoRA, dual LoRA with fixed fusion, and dual LoRA with a learned fusion gate. On the recorded ROF protocol averaged over three training seeds, periocular LoRA gives the highest masked ROC-AUC (0.8989$\pm$0.008), whereas learned gating gives the highest sunglasses ROC-AUC (0.8998$\pm$0.005) and neutral AUC (0.9441$\pm$0.002). A normalized exact-name audit found one ROF training identity also present in MFR2, affecting 18 of 820 pairs; MFR2 multi-seed results (mean AUC 0.8405--0.9554) are reported with this contamination caveat and are not interpreted as fully unseen-identity generalization. On RMFRD (exploratory, with duplicate-path and quality-metadata limitations), fixed fusion achieves the highest mean AUC (0.7184$\pm$0.011). The evidence therefore supports a conditional conclusion: parameter-efficient adaptation and periocular fusion improve selected operating points, but learned gating is not uniformly dominant across conditions, datasets, or false-accept regimes.
\end{abstract}
\begin{IEEEkeywords}
Face verification, occlusion robustness, masked-face verification, periocular representation, DINOv2, LoRA, learned fusion, cross-dataset evaluation.
\end{IEEEkeywords}

\noindent\textit{Disclosure statement.} Generative AI tools (OpenAI ChatGPT)
were used during preparation of this manuscript to assist with drafting and
revising portions of the prose, language editing, restructuring of sections, and
consistency checking across the manuscript and the accompanying artifact. The
research design, experimental execution, data preparation, model training,
evaluation, and all reported numerical results were produced by the authors'
own workflow; no generative AI tool generated, modified, or selected any
experimental result, and every reported value was recomputed from the archived
raw scores by the released verification script. The authors reviewed and take
full responsibility for the content of this publication.

\section{Introduction}
Face verification determines whether two images belong to the same identity. This differs from closed-set recognition, in which the system assigns an image to one of the identities observed during training. Verification is especially demanding under masks and sunglasses because the visible evidence changes across conditions, and because a false accept at a low false-accept rate can be more consequential than a modest loss in average ranking quality.

A whole-face representation is a natural baseline, but it is not uniformly reliable under occlusion. A mask removes much of the lower face, while sunglasses alter the eye region that is especially important for periocular recognition. A periocular crop can retain visible identity information under masks, but it can also discard complementary cues present in the full face. These observations motivate a controlled question: can a parameter-efficient system learn when to rely on the whole-face branch and when to rely on the periocular branch?

This paper investigates that question with a DINOv2 backbone and LoRA adaptation. The central methodological design is not simply the combination of established components. Instead, the study tests whether adaptive fusion earns its place through a five-configuration ablation and through cross-condition and cross-dataset verification protocols. The evaluation uses person-level identity-disjoint splits, saved raw scores, and explicit reporting of negative findings. MFR2 is evaluated as an identity-overlap diagnostic (see Section V-B), not as a clean cross-identity benchmark.

The paper makes five contributions. First, it provides a controlled comparison of frozen whole-face, whole-face LoRA, periocular LoRA, fixed dual-branch fusion, and learned dual-branch gating. Second, it normalizes condition-specific ROF directory identities before creating person-disjoint partitions, preventing the same individual from being represented as different identities across conditions. Third, it evaluates neutral enrollment against masked and sunglasses probes with ROC-AUC, EER, TAR at low FAR, score distributions, and confidence intervals. Fourth, it reports an external MFR2 evaluation as an identity-overlap diagnostic (not a clean cross-identity benchmark). Fifth, it documents an RMFRD audit limitation rather than using the external dataset to support an unqualified generalization claim.

\section{Related Work}
\label{sec:related}
\subsection{Face verification}
Deep face verification systems learn embeddings in which same-identity images are close and different-identity images are separated. Cosine similarity remains a common comparison rule because it is simple, scale-insensitive after normalization, and directly compatible with ROC analysis. Margin-based softmax losses have driven steady improvement on clean-face benchmarks. SphereFace introduced the angular margin concept \cite{sphereface}, CosFace replaced it with an additive cosine margin \cite{cosface}, ArcFace added an angular margin on the hypersphere \cite{arcface}, and AdaFace made the margin adaptive to sample quality \cite{adaface}. MagFace jointly learns embeddings and a quality-aware margin from a single network \cite{magface}. In this work, classification over training identities is treated only as an optimization device; all primary claims use verification scores on identities withheld from training.

\subsection{Masked and occluded face recognition}
Masked-face recognition became a distinct subfield during the COVID-19 pandemic. Anwar and Raychowdhury introduced the MaskTheFace simulator and the MFR2 real-mask dataset \cite{mfr2}. Ge \textit{et al.} built the MAFA dataset and LLE-CNNs for masked-face detection \cite{llecnn}. Wang \textit{et al.} constructed the RMFRD family of datasets (MFDD, RMFRD, SMFRD) \cite{rmfrd}. A comprehensive survey by Zeng, Veldhuis, and Spreeuwers covers occlusion-robust face recognition up to 2021 \cite{surveyocc}. More recent methods include the end-to-end FROM approach that dynamically masks corrupted features \cite{from2022}, the pairwise differential Siamese mask-learning of Song \textit{et al.} \cite{masklearn}, and the self-restrained triplet loss of Boutros \textit{et al.} that unmasks face embeddings \cite{umr}. Neto \textit{et al.} proposed FocusFace (multi-task contrastive learning) \cite{focusface} and analyzed the generalization of masked-face models to broader occlusion \cite{beyondmasks}. The MFR 2021 competition at IJCB and the InsightFace track report provide strong benchmarks on real masked data \cite{mfr2021,insightfacemfr}. Hsu \textit{et al.} explored synthesis-to-reality transfer \cite{hsu2022}. The ROF dataset of Erakın, Demir, and Ekenel offers real-world upper-face (sunglasses) and lower-face (mask) occlusions and dedicated evaluation protocols \cite{rof}. Recent work also studies face--periocular complementarity directly: Flexible Biometrics Recognition aligns face and periocular representations with attention and prompt tuning \cite{fbr2024}. This is especially relevant to our setting because it treats the two regions as complementary rather than assuming that either is uniformly reliable.

\subsection{Periocular recognition}
Periocular biometrics uses the eye region and surrounding upper-face context. Park, Jillela, Ross, and Jain demonstrated feasibility in the visible spectrum, showing rank-one accuracy of 87.3\% on FRGC v2.0 and confirming that the periocular region carries much of the identity signal \cite{periocular}. Earlier work by the same group established the feasibility study at BTAS 2009. More recent periocular research transfers identity information from full faces to periocular embeddings through consistent knowledge distillation \cite{ckd2024}, illustrating that the two representations can be complementary even when only the periocular region is available at inference. Periocular recognition is therefore a plausible fallback under lower-face occlusion, but sunglasses remain a direct challenge because they obscure the region on which the representation depends.

\subsection{Self-supervised visual representations}
DINOv2 provides transferable features from large-scale self-supervised training on 142M images \cite{dinov2}. Its predecessor DINO emerged properties in vision transformers \cite{dino}. Masked autoencoders (MAE) showed that reconstructing masked patches yields scalable vision learners \cite{mae}. Contrastive methods such as SimCLR \cite{simclr} and CLIP \cite{clip} demonstrated that self-supervision on massive image-text or image-image pairs produces strong linear-probe features. These representations are attractive for our setting because the occlusion datasets are modest and a strong visual prior reduces the adaptation burden.

\subsection{Parameter-efficient adaptation}
LoRA introduces low-rank trainable updates while keeping the original network frozen \cite{lora}, reducing trainable parameters by orders of magnitude. Prompt-based methods such as visual prompt tuning (VPT) \cite{vpt} and the power-of-scale prompt tuning \cite{prompt} prepend learnable tokens or prefixes. Adapter-style modules such as AdaptFormer insert lightweight bottlenecks while leaving the pretrained transformer largely unchanged \cite{adaptformer}. We use LoRA because it directly modifies the attention projections and integrates cleanly with the DINOv2 backbone.

\subsection{Score fusion and quality-aware gating}
Biometric fusion has a long history; Ross and Jain surveyed score- and decision-level strategies \cite{fusionRossJain}. Quality-dependent fusion weights samples by their estimated reliability \cite{quality}. Learned gating mechanisms condition the fusion weight on sample-specific cues. Our gate consumes branch embeddings and crop-quality descriptors to produce a scalar mixing weight, distinct from margin-adaptive losses \cite{magface,adaface} because it operates at the embedding combination stage rather than during classifier training.

\section{Research Questions and Hypotheses}
The study addresses three questions. RQ1 asks whether LoRA adaptation improves verification over a frozen whole-face DINOv2 representation. RQ2 asks whether a periocular branch improves masked verification while preserving performance under sunglasses. RQ3 asks whether learned fusion improves over a fixed 0.5/0.5 combination and whether the improvement transfers to MFR2 and RMFRD.

The corresponding hypotheses are deliberately falsifiable. H1 predicts that whole-face LoRA improves over frozen whole-face DINOv2. H2 predicts that periocular LoRA improves masked verification but may not improve sunglasses verification. H3 predicts that learned gating improves at least some low-FAR operating points relative to fixed fusion. The results do not require every hypothesis to hold on every dataset.

\section{Method}
\subsection{Representation branches}
Let $x_f$ be a standardized whole-face crop and $x_p$ be a periocular crop containing both eyes, eyebrows, the upper nose, and local context. Let $E_f$ and $E_p$ denote DINOv2 encoders and let $P_f$ and $P_p$ denote trainable projection heads. The normalized branch embeddings are
\begin{equation}
 z_f=\operatorname{norm}(P_f(E_f(x_f))),\qquad z_p=\operatorname{norm}(P_p(E_p(x_p))).
\end{equation}
The whole-face and periocular branches use the same image resolution and corresponding image normalization. Landmark failures are recorded in the crop manifest and are not silently replaced by a whole-face image.

\subsection{Fixed and learned fusion}
The fixed-fusion baseline combines both branches with equal weight:
\begin{equation}
 z_{\mathrm{fixed}}=\operatorname{norm}(0.5z_f+0.5z_p).
\end{equation}
The learned-gating system computes a scalar reliability value
\begin{equation}
 g=\sigma\left(\operatorname{MLP}([z_f;z_p;q_f;q_p])\right),
\end{equation}
where $q_f$ and $q_p$ are available crop-quality descriptors. The final embedding is
\begin{equation}
 z_{\mathrm{gate}}=\operatorname{norm}(gz_f+(1-g)z_p).
\end{equation}
The gate is treated as a learned fusion mechanism rather than a calibrated reliability estimator. Its input contains both branch embeddings and four auxiliary descriptors, but the available quality fields are sparse or constant in parts of the evaluation data. Consequently, this paper does not interpret the scalar gate value as a calibrated quality score. Direct gate-distribution analysis and an embedding-only gate ablation remain useful follow-up experiments.

\subsection{LoRA adaptation and training objective}
LoRA updates are inserted into the selected attention projection layers. The encoder's original parameters remain frozen, while LoRA parameters and projection/fusion heads are optimized. Let $h_i$ be the embedding for a training image and $y_i$ its training identity. The primary training objective is identity cross-entropy over the training identities:
\begin{equation}
 \mathcal{L}_{\mathrm{ID}}=-\sum_i\log p(y_i\mid h_i).
\end{equation}
No test identity is used to define classifier classes, tune the architecture, or select thresholds. The verification score for two normalized embeddings is
\begin{equation}
 s(i,j)=z_i^\top z_j.
\end{equation}

\subsection{Ablation configurations}
The five configurations in Table~\ref{tab:ablations} cover the frozen baseline, the adaptation effect, the periocular-only effect, fusion without gating, and the proposed learned gate.
\begin{table}[t]
\centering
\caption{Five required configurations.}
\label{tab:ablations}
\begin{tabular}{cll}
\toprule
ID & Configuration & Purpose\\
\midrule
M0 & Frozen whole-face DINOv2 & Frozen baseline\\
M1 & Whole-face DINOv2 + LoRA & Adaptation effect\\
M2 & Periocular DINOv2 + LoRA & Periocular-only effect\\
M3 & Dual LoRA + fixed fusion & Fusion without gating\\
M4 & Dual LoRA + learned gate & Proposed adaptive fusion\\
\bottomrule
\end{tabular}
\end{table}

\section{Datasets and Protocol}
\subsection{ROF identity audit}
ROF is the primary dataset. The audit identified 7,752 images and 400 condition-specific directory identities before normalization. Since condition-specific names can represent the same person, names ending in condition markers such as ``wearing mask'' and ``wearing sunglasses'' were normalized before partitioning. The final person-disjoint complete-case manifest contained approximately 128 training identities, 26 validation identities, and 28 test identities in the recorded evaluation. Exact counts and the final manifest hash are part of the accompanying artifact.

Training uses only training identities. Validation identities are reserved for development decisions and threshold selection. Test identities are used once for final verification reporting. ROF evaluation uses a neutral enrollment image and probes from either the masked or sunglasses condition. All five systems use the same saved pair lists.

\subsection{MFR2 official pairs}
MFR2 is used as an external masked-face verification benchmark. Models are trained on ROF and evaluated on MFR2 without target-domain fine-tuning. The complete-case protocol contains 820 official pairs: 412 genuine and 408 impostor pairs. Results are reported separately from ROF because the pair construction and image source differ. A normalized exact-name audit found a single ROF training identity also present in MFR2, which affects 18 of the 820 evaluated pairs; MFR2 is therefore not fully identity-disjoint from the ROF training set, but the contamination is limited. The zero-overlap statement (ROF$\cap$MFR2$=\varnothing$) in earlier drafts was incorrect and has been removed.

\subsection{RMFRD exploratory transfer}
RMFRD is included only as exploratory cross-dataset transfer. The recorded v2 evaluation contains 22,554 pairs, including 1,074 genuine and 21,480 impostor comparisons. The audit reports exact duplicates across nominal identities and incomplete quality fields; these issues can distort impostor distributions and prevent a clean independent-identity claim. RMFRD is therefore not used as primary evidence until the audit is resolved. The recorded RMFRD manifest also carries blank quality descriptors, so the gate received placeholder quality inputs there; for this reason the learned-gate result on RMFRD should not be read as evidence about gate behavior under valid quality cues.

\subsection{Metrics and uncertainty}
We report ROC-AUC, EER, and TAR at FAR $10^{-2}$ and $10^{-3}$. Genuine and impostor scores are saved as CSV files. Bootstrap intervals are reported with the metrics where available. Because multiple pairs can share an identity, the preferred final uncertainty analysis is identity-aware paired bootstrap: identities are resampled, all associated pairs are included, and model differences are computed on the same resampled identity sets.

\subsection{Protocol integrity and audit controls}
The protocol audit is part of the experimental method rather than a post-hoc presentation step. Three checks are particularly important. First, ROF identities are normalized before splitting so that condition-specific directory names cannot place the same person in both training and test partitions. Second, every reported verification protocol is materialized as a frozen pair list; all five variants are evaluated on exactly the same pairs, which makes paired model comparisons meaningful. Third, external datasets are not assumed to be independent merely because their files originate from a different corpus. Identity names and image paths are audited explicitly.

This audit changed the interpretation of two external evaluations. MFR2 shares identities with the ROF training set, so it is retained as a stress test but not used as evidence of unseen-identity transfer. In addition, the first multi-seed MFR2 evaluator incorrectly cached embeddings by identity rather than by image path. That implementation collapsed multiple images of one person and produced artificial self-comparisons. The final evaluator caches by the concrete $(\texttt{face\_path},\texttt{peri\_path})$ tuple, rejects exact image self-pairs, and records the cache convention in the metric metadata. The corrected MFR2 values in this paper are generated only by that image-level evaluator. RMFRD, separately, contains duplicate-path and incomplete-quality issues and is therefore marked exploratory.

\begin{table}[t]
\centering
\caption{Evaluation protocols and interpretation after audit.}
\label{tab:protocols}
\resizebox{\columnwidth}{!}{%
\begin{tabular}{lrrrl}
\toprule
Protocol & Genuine & Impostor & IDs & Interpretation\\
\midrule
ROF masked & 58 & 464 & 9 & Primary, person-disjoint\\
ROF sunglasses & 327 & 7,521 & 24 & Primary, person-disjoint\\
ROF neutral & 677 & 18,279 & 28 & In-domain control\\
MFR2 & 412 & 408 & 53 & Overlap stress test\\
RMFRD v2 & 1,074 & 21,480 & 341 & Exploratory\\
\bottomrule
\end{tabular}%
}
\end{table}

Table~\ref{tab:protocols} also makes the low-FAR resolution explicit. With $N_I$ impostor comparisons, the smallest non-zero empirical FAR increment is $1/N_I$. For ROF masked, $1/464\approx2.16\times10^{-3}$, so TAR at FAR $10^{-3}$ lies below the protocol's empirical resolution. We retain it only for completeness and do not use that operating point as the primary basis for a superiority claim.

\section{Experimental Settings}
All configurations share one frozen evaluation protocol and the same optimization budget; only the component named in the ablation changes. Measured settings from the run logs are as follows.

\textbf{Architecture.} Backbone: \texttt{facebook/dinov2-base} (DINOv2 ViT-B/14, embedding dimension 768). Inputs: $224\times224$ RGB crops with ImageNet mean/std normalization. Projection heads: linear $768\rightarrow768$. Gate MLP: linear $1540\rightarrow128$, GELU, linear $128\rightarrow1$, sigmoid output. LoRA: rank $r=8$, scaling $\alpha=16$, dropout $0.1$, applied to the \texttt{query}, \texttt{key}, and \texttt{value} projections of every attention layer; encoder base weights remain frozen.

\textbf{Optimization.} Identity cross-entropy over the 128 training identities; AdamW, learning rate $10^{-4}$, weight decay $10^{-4}$, batch size 16, source-weighted sampling. Final models are trained independently with seeds 7, 42, and 123. Measured throughput: 325 batches/epoch over the 5,194-image training split, five epochs. No test identity defines a classifier class, tunes the architecture, or selects thresholds.

\subsection{Statistical reporting protocol}
The final report separates three sources of evidence. First, each training seed produces a complete checkpoint and an independent set of verification scores. We summarize these runs with the arithmetic mean and sample standard deviation across seeds $\{7,42,123\}$. Second, for a fixed checkpoint we estimate sampling uncertainty by resampling identities rather than individual pair rows. This preserves within-identity dependence among comparisons. Third, when two models are compared, the same resampled identities are used for both systems and the bootstrap distribution is formed from the metric difference, rather than from two independently bootstrapped confidence intervals. This paired construction is the relevant uncertainty estimate for claims such as gated versus fixed fusion.

Let $m(A;\mathcal{S})$ denote a verification metric for model $A$ on all pairs associated with a resampled identity multiset $\mathcal{S}$. For models $A$ and $B$, each bootstrap replicate is
\begin{equation}
 \Delta_b=m(A;\mathcal{S}_b)-m(B;\mathcal{S}_b),
\end{equation}
and the reported 95\% interval is given by the 2.5th and 97.5th percentiles of $\{\Delta_b\}$. This design avoids treating highly correlated pair rows as independent observations. The bootstrap RNG seed is metadata for the uncertainty calculation and is recorded separately from the training seed that generated the checkpoint.

The paper does not perform significance testing over the three training seeds themselves. With only three initializations, a parametric test on seed means would be poorly determined and could create a false sense of precision. Instead, seed-to-seed variation is reported descriptively, while identity-aware paired bootstrap intervals quantify protocol sampling uncertainty. Claims in the Discussion distinguish these two quantities explicitly.

\subsection{Parameter efficiency}
\begin{table}[h]
\centering
\caption{Trainable-parameter accounting measured at model construction.}
\label{tab:params}
\begin{tabular}{lrrr}
\toprule
ID & Total parameters & Trainable & Trainable \%\\
\midrule
M0 & 87\,171\,072 & 590\,592 & 0.68\\
M1 & 87\,613\,440 & 1\,032\,960 & 1.18\\
M2 & 175\,226\,880 & 2\,065\,920 & 1.18\\
M3 & 175\,226\,880 & 2\,065\,920 & 1.18\\
M4 & 175\,424\,257 & 2\,263\,297 & 1.29\\
\bottomrule
\end{tabular}
\end{table}

Table~\ref{tab:params} reports the trainable-parameter accounting measured at model construction. Every adapted variant therefore trains at most 2.26~M parameters (1.29\%) of a 175~M-parameter dual-encoder system. Hardware, runtime, peak memory, software versions, and checkpoint SHA-256 values are stored in the artifact rather than inferred from scores.

\subsection{Reproducibility safeguards}
Several implementation safeguards were added during the audit and are retained in the final code. They are summarized in Table~\ref{tab:safeguards}. These checks matter because an apparently strong verification result can be produced by bookkeeping errors even when the model itself is unchanged. In particular, the MFR2 cache-key defect demonstrated that image identity and person identity must not be conflated when one person contributes multiple images.

\begin{table}[t]
\centering
\caption{Safeguards used for the final reported experiments.}
\label{tab:safeguards}
\begin{tabular}{p{0.24\columnwidth}p{0.64\columnwidth}}
\toprule
Check & Final handling\\
\midrule
Person leakage & Normalize ROF identities before person-disjoint split\\
Pair consistency & Reuse frozen pair CSVs for every variant and seed\\
MFR2 cache key & Key embeddings by concrete face/periocular paths\\
Exact self-pairs & Raise an evaluation error rather than score silently\\
Result provenance & Store checkpoint and pair-list SHA-256 hashes\\
Seed provenance & Separate training seed from bootstrap RNG seed\\
Smoke tests & Exclude smoke/partial checkpoints from paper tables\\
External quality & Mark missing/placeholder descriptors explicitly\\
\bottomrule
\end{tabular}
\end{table}

\section{Results}
\subsection{ROF masked probes}
Table~\ref{tab:rofmasked} reports the ROF masked verification results averaged over three training seeds (7, 42, 123), with the corresponding ROC curves in Fig.~\ref{fig:rofmaskedroc}. Periocular LoRA produces the highest mean AUC (0.8989), while fixed fusion gives the lowest mean EER (0.1695) and the highest recorded TAR at both reported low-FAR points (0.2874 at FAR $10^{-2}$ and 0.2069 at FAR $10^{-3}$). This distinction is important: the ranking winner and the low-FAR operating-point winner are different systems. Learned gating is close in AUC (0.8942) and at FAR $10^{-2}$ (0.2644), but does not produce the highest AUC, lowest EER, or highest TAR. Fixed fusion provides the second-best masked AUC (0.8949). Whole-face LoRA improves substantially over the frozen baseline (mean AUC +0.0315). The standard deviations across seeds are small (max 0.013 for AUC, max 0.016 for EER), indicating stable training dynamics.

\begin{table}[t]
\centering
\caption{ROF masked probes, neutral enrollment. Values are mean AUC, mean EER, mean TAR at FAR $10^{-2}$, and mean TAR at FAR $10^{-3}$ over three seeds (7, 42, 123); subscripts show standard deviation across seeds.}
\label{tab:rofmasked}
\resizebox{\columnwidth}{!}{%
\begin{tabular}{lrrrr}
\toprule
Variant & AUC & EER & TAR$_{.01}$ & TAR$_{.001}$\\
\midrule
Frozen whole & .8306$_{\pm.013}$ & .2759$_{\pm.013}$ & .2011$_{\pm.040}$ & .0000$_{\pm.000}$\\
Whole LoRA & .8621$_{\pm.006}$ & .2011$_{\pm.016}$ & .2701$_{\pm.061}$ & .1264$_{\pm.050}$\\
Periocular LoRA & \textbf{.8989}$_{\pm.008}$ & .1803$_{\pm.012}$ & .1207$_{\pm.079}$ & .0402$_{\pm.026}$\\
Fixed fusion & .8949$_{\pm.009}$ & \textbf{.1695}$_{\pm.013}$ & \textbf{.2874}$_{\pm.100}$ & \textbf{.2069}$_{\pm.130}$\\
Learned gating & .8942$_{\pm.005}$ & .1853$_{\pm.011}$ & .2644$_{\pm.078}$ & .1552$_{\pm.119}$\\
\bottomrule
\end{tabular}%
}

{\footnotesize 95\% AUC intervals (probe-identity bootstrap, seed 42): Frozen whole 0.727--0.944; Whole LoRA 0.754--0.950; Periocular LoRA 0.833--0.953; Fixed fusion 0.749--0.976; Learned gating 0.807--0.961. TAR at FAR $10^{-3}$ is below the FAR resolution of this protocol (${\approx}2.2\times10^{-3}$ given 464 impostor pairs) and is reported for completeness only.}
\end{table}

\begin{figure}[t]
\centering
\includegraphics[width=\columnwidth]{fig_rof_masked_roc.png}
\caption{ROF masked verification ROC curves (seed 42). The small number of eligible masked test identities makes low-FAR estimates discrete and uncertain.}
\label{fig:rofmaskedroc}
\end{figure}

\subsection{ROF sunglasses probes}
Table~\ref{tab:rofsunglasses} shows the sunglasses results averaged over three seeds (Fig.~\ref{fig:rofsunglassesroc}). Learned gating provides the highest mean AUC (0.8998) and lowest mean EER (0.1528), while fixed fusion provides competitive AUC (0.8933) and whole-face LoRA is close (0.8961). Periocular LoRA is substantially weaker (mean AUC 0.7919), supporting the hypothesis that sunglasses can directly impair a periocular-only representation. The standard deviations across seeds remain small (max 0.008 for AUC, max 0.010 for EER), confirming stable rankings.

\begin{table}[t]
\centering
\caption{ROF sunglasses probes, neutral enrollment. Mean values and standard deviations over three seeds (7, 42, 123).}
\label{tab:rofsunglasses}
\resizebox{\columnwidth}{!}{%
\begin{tabular}{lrrrr}
\toprule
Variant & AUC & EER & TAR$_{.01}$ & TAR$_{.001}$\\
\midrule
Frozen whole & .8754$_{\pm.003}$ & .1753$_{\pm.005}$ & .1091$_{\pm.008}$ & .0194$_{\pm.002}$\\
Whole LoRA & .8961$_{\pm.004}$ & .1609$_{\pm.007}$ & \textbf{.1590}$_{\pm.038}$ & .0285$_{\pm.009}$\\
Periocular LoRA & .7919$_{\pm.007}$ & .2675$_{\pm.010}$ & .0571$_{\pm.011}$ & .0061$_{\pm.003}$\\
Fixed fusion & .8933$_{\pm.008}$ & .1579$_{\pm.008}$ & .1407$_{\pm.006}$ & .0204$_{\pm.005}$\\
Learned gating & \textbf{.8998}$_{\pm.005}$ & \textbf{.1528}$_{\pm.008}$ & .1570$_{\pm.033}$ & \textbf{.0316}$_{\pm.020}$\\
\bottomrule
\end{tabular}%
}

{\footnotesize 95\% AUC intervals (probe-identity bootstrap, seed 42): Frozen whole 0.820--0.913; Whole LoRA 0.860--0.924; Periocular LoRA 0.693--0.853; Fixed fusion 0.862--0.924; Learned gating 0.875--0.933.}
\end{table}
\begin{figure}[t]
\centering
\includegraphics[width=\columnwidth]{fig_rof_sunglasses_roc.png}
\caption{ROF sunglasses verification ROC curves (seed 42). Fusion reduces the degradation observed for periocular-only LoRA.}
\label{fig:rofsunglassesroc}
\end{figure}

\subsection{MFR2 external evaluation}
MFR2 provides an external masked-face verification benchmark. Models are trained on ROF and evaluated on MFR2 without target-domain fine-tuning. The complete-case protocol contains 820 official pairs: 412 genuine and 408 impostor pairs. A normalized exact-name audit found a single ROF training identity also present in MFR2, which affects 18 of the 820 evaluated pairs. MFR2 is therefore reported as a limited-overlap diagnostic; its results are reported with this contamination caveat and are not interpreted as fully unseen-identity generalization evidence. Results are reported separately from ROF because the pair construction and image source differ.

Table~\ref{tab:mfr2} reports the mean and standard deviation over three training seeds (7, 42, 123), with the seed-42 ROC curves in Fig.~\ref{fig:mfr2roc}. Learned gating achieves the highest mean AUC (0.9554) and lowest mean EER (0.1148). Periocular LoRA gives the highest mean TAR at FAR $10^{-3}$ (0.1262), narrowly above fixed fusion (0.1246). Whole-face LoRA and fixed fusion are competitive. The standard deviations across seeds are small (max 0.012 for AUC), indicating stable rankings.

\begin{table}[t]
\centering
\caption{MFR2 official complete-case pairs: 412 genuine and 408 impostor pairs. Values are mean and standard deviation over three seeds (7, 42, 123). One ROF training identity also appears in MFR2 (18 of 820 pairs); the protocol is therefore not strictly identity-disjoint.}
\label{tab:mfr2}
\resizebox{\columnwidth}{!}{%
\begin{tabular}{lrrrr}
\toprule
Variant & AUC & EER & TAR$_{.01}$ & TAR$_{.001}$\\
\midrule
Frozen whole & .8405$_{\pm.003}$ & .2427$_{\pm.002}$ & .1529$_{\pm.021}$ & .0032$_{\pm.001}$\\
Whole LoRA & .9397$_{\pm.005}$ & .1380$_{\pm.007}$ & .3382$_{\pm.061}$ & .0663$_{\pm.059}$\\
Periocular LoRA & .9001$_{\pm.012}$ & .1824$_{\pm.018}$ & .2953$_{\pm.041}$ & \textbf{.1262}$_{\pm.051}$\\
Fixed fusion & .9381$_{\pm.012}$ & .1495$_{\pm.019}$ & .3673$_{\pm.056}$ & .1246$_{\pm.021}$\\
Learned gating & \textbf{.9554}$_{\pm.002}$ & \textbf{.1148}$_{\pm.005}$ & \textbf{.4636}$_{\pm.029}$ & .0089$_{\pm.001}$\\
\bottomrule
\end{tabular}%
}

{\footnotesize One ROF training identity also appears in MFR2, affecting 18 of 820 pairs (2.2\%); the protocol is not strictly identity-disjoint, and these values are not interpreted as fully unseen-identity generalization evidence.}
\end{table}
\begin{figure}[t]
\centering
\includegraphics[width=\columnwidth]{fig_mfr2_roc.png}
\caption{MFR2 official-pair ROC curves (seed 42). One ROF training identity also appears in MFR2 (18 of 820 pairs).}
\label{fig:mfr2roc}
\end{figure}

\subsection{RMFRD exploratory transfer}
RMFRD v2 is included only as exploratory cross-dataset transfer. The recorded v2 evaluation contains 22,554 pairs, including 1,074 genuine and 21,480 impostor comparisons. The audit reports exact duplicates across nominal identities and incomplete quality fields; these issues can distort impostor distributions and prevent a clean independent-identity claim. RMFRD is therefore not used as primary evidence until the audit is resolved. The recorded RMFRD manifest also carries blank quality descriptors, so the gate received placeholder quality inputs there; for this reason the learned-gate result on RMFRD should not be read as evidence about gate behavior under valid quality cues.

Table~\ref{tab:rmfrd} reports the mean and standard deviation over three training seeds (7, 42, 123), and Fig.~\ref{fig:crossauc} compares mean AUC across all protocols. Fixed fusion gives the highest mean AUC (0.7184), followed by periocular LoRA (0.7144) and whole-face LoRA (0.7084). Learned gating gives 0.7063. The standard deviations are small (max 0.014), indicating stable rankings. Because the audit reports cross-identity exact duplicate paths, these results are not used to establish independent-dataset superiority.

\begin{table}[t]
\centering
\caption{RMFRD v2 exploratory transfer, mean$\pm$std over three seeds (7, 42, 123). Interpret with the duplicate-path and quality-metadata audit limitation.}
\label{tab:rmfrd}
\resizebox{\columnwidth}{!}{%
\begin{tabular}{lrrrr}
\toprule
Variant & AUC & EER & TAR$_{.01}$ & TAR$_{.001}$\\
\midrule
Frozen whole & .6955$_{\pm.003}$ & .3517$_{\pm.001}$ & .0239$_{\pm.003}$ & .0047$_{\pm.000}$\\
Whole LoRA & .7084$_{\pm.014}$ & .3483$_{\pm.015}$ & .0410$_{\pm.008}$ & .0078$_{\pm.002}$\\
Periocular LoRA & .7144$_{\pm.007}$ & .3410$_{\pm.006}$ & .0323$_{\pm.003}$ & .0050$_{\pm.001}$\\
Fixed fusion & \textbf{.7184}$_{\pm.011}$ & \textbf{.3392}$_{\pm.007}$ & \textbf{.0438}$_{\pm.007}$ & .0065$_{\pm.001}$\\
Learned gating & .7063$_{\pm.006}$ & .3436$_{\pm.007}$ & .0341$_{\pm.002}$ & .0068$_{\pm.002}$\\
\bottomrule
\end{tabular}%
}

{\footnotesize 95\% AUC intervals (probe-identity bootstrap, seed 42): Frozen whole 0.674--0.721; Whole LoRA 0.699--0.745; Periocular LoRA 0.694--0.742; Fixed fusion 0.695--0.742; Learned gating 0.684--0.733. RMFRD is exploratory because of the duplicate-path and quality-metadata audit findings.}
\end{table}
\begin{figure}[t]
\centering
\includegraphics[width=\columnwidth]{fig_auc_cross_dataset.png}
\caption{Mean ROC-AUC over three training seeds across the evaluated protocols. Error bars show one standard deviation. One ROF training identity also appears in MFR2 (18 of 820 pairs); RMFRD is exploratory because of the audit findings.}
\label{fig:crossauc}
\end{figure}

\subsection{Score distributions and qualitative error analysis}
Figure~\ref{fig:scores} shows the gated-fusion genuine and impostor score distributions for ROF masked probes. Overlap between the distributions explains why a high AUC does not automatically imply a high TAR at a stringent FAR. The final artifact retains raw scores so that false accepts and false rejects can be inspected by identity, condition, crop status, and quality descriptor.
\begin{figure}[t]
\centering
\includegraphics[width=\columnwidth]{fig_rof_masked_scores.png}
\caption{Gated-fusion score distributions on ROF masked probes.}
\label{fig:scores}
\end{figure}

\subsection{Reading AUC and low-FAR metrics together}
AUC, EER, and TAR at a fixed FAR answer different questions. AUC summarizes ranking over the full score range, EER identifies the point where false-accept and false-reject rates are equal, and TAR at a prescribed FAR focuses on a security-sensitive operating region. The experiments illustrate why these metrics should not be collapsed into a single winner. On ROF masked probes, periocular LoRA has the highest mean AUC, while fixed fusion has the lowest mean EER and the highest recorded TAR at the two reported FAR targets. On ROF sunglasses, learned gating leads in AUC and EER, while whole-face LoRA narrowly leads at FAR $10^{-2}$ (0.1590 vs. 0.1570 for learned gating) and learned gating leads at FAR $10^{-3}$. The method comparison is therefore operating-point dependent.

The three-seed standard deviations describe sensitivity to training initialization, while the bootstrap intervals describe sampling uncertainty under a fixed trained model/protocol. These are complementary quantities and are not interchangeable. We consequently avoid treating a small mean difference as statistically decisive when paired bootstrap intervals include zero. This distinction is particularly important for the gated-versus-fixed comparison on sunglasses, where the mean AUC difference is positive but modest.

\section{Discussion}
The results answer the research questions conditionally. LoRA generally improves over frozen whole-face DINOv2 on ROF ranking measures, supporting H1. Periocular LoRA is strongest on the recorded ROF masked AUC (mean 0.8989 vs. 0.8949 for fixed fusion and 0.8942 for learned gating) but weak under sunglasses (mean 0.7919), supporting H2 in a condition-dependent form. Learned gating improves selected measures on ROF sunglasses (mean AUC 0.8998 vs. 0.8933 for fixed fusion), but the fixed-fusion and whole/periocular baselines remain competitive at other operating points and on RMFRD, so H3 is only partially supported.

The main scientific finding is therefore not that learned gating is universally best. It is that branch reliability is condition- and dataset-dependent. Under a mask, the periocular branch can provide useful evidence, but under sunglasses it is directly impaired. A fixed fusion can be robust because it does not overreact to noisy branch estimates. A learned gate can improve ranking when its training distribution matches the test condition, but external transfer can expose a mismatch between learned quality cues and target-domain reliability.

On the primary ROF sunglasses comparison, the gated-fusion vs. fixed-fusion AUC difference across three seeds is $+0.0065\pm0.0105$ (mean $\pm$ std). One seed (123) shows a positive identity-aware CI $[+0.0046, +0.0254]$, while the other two cross zero. The multi-seed mean is therefore positive but not statistically clean; the gating advantage on ROF sunglasses remains modest. By contrast, fixed fusion over frozen whole-face gives a large mean AUC gain of $+0.0643\pm0.0156$ on masked probes. The corresponding gains are $+0.0179\pm0.0102$ on sunglasses and $+0.0111\pm0.0041$ on neutral, with positive point estimates across seeds. The per-seed identity-aware intervals, however, do not exclude zero in every case. On the neutral in-domain control, learned gating also attains the highest mean AUC (0.9441$\pm$0.002) and lowest mean EER (0.1101$\pm$0.002), with fixed fusion close (0.9424$\pm$0.002) and periocular-only LoRA markedly weaker (0.8832$\pm$0.006).

On MFR2, the multi-seed mean AUCs range from 0.8405 (frozen whole) to 0.9554 (gated fusion), with small standard deviations (max 0.012). Learned gating achieves the highest mean AUC (0.9554$\pm$0.002) and lowest mean EER (0.1148$\pm$0.005). However, the MFR2 audit found one ROF training identity also present in MFR2 (affecting 18 of 820 pairs, 2.2\%), so the protocol is not strictly identity-disjoint and these results are not interpreted as fully unseen-identity generalization. On RMFRD (exploratory, with audit limitations), fixed fusion achieves the highest mean AUC (0.7184$\pm$0.011), followed by periocular LoRA (0.7144$\pm$0.007) and whole-face LoRA (0.7084$\pm$0.014); learned gating gives 0.7063$\pm$0.006. The RMFRD audit reports exact duplicate paths across nominal identities, so these results are exploratory only.

\subsection{Positioning relative to related work}
Direct numerical ranking against prior masked-face systems is difficult because published studies use different training identities, mask synthesis procedures, pair lists, and target-domain adaptation policies. We therefore compare design choices rather than claim a protocol-mismatched state-of-the-art result. ROF itself was introduced to expose the gap between synthetic and real-world occlusion \cite{rof}; our results reinforce that motivation because performance changes sharply with the location of the occluder. MFR 2021 and the InsightFace challenge similarly demonstrate the need for masked-face-specific evaluation at scale \cite{mfr2021,insightfacemfr}.

Existing occlusion-robust methods attack the problem in several ways. FROM learns masks over corrupted features \cite{from2022}; pairwise differential Siamese learning explicitly models occlusion differences \cite{masklearn}; self-restrained triplet learning modifies the embedding objective for masked faces \cite{umr}; and FocusFace uses multi-task contrastive learning \cite{focusface}. Our design instead asks how far a frozen self-supervised visual backbone can be adapted with a small trainable parameter budget and an explicit face/periocular decomposition. This parameter-efficient emphasis is closer to LoRA/VPT/adapter transfer \cite{lora,vpt,adaptformer} than to full masked-face retraining.

The face--periocular decomposition is also consistent with newer multimodal biometric work. Flexible Biometrics Recognition uses attention alignment and prompt tuning to bridge face and periocular modalities \cite{fbr2024}, while consistent knowledge distillation transfers face identity information into a periocular embedding \cite{ckd2024}. Our experiments address a different question: rather than learning a universal cross-modal matcher, we measure when whole-face, periocular, fixed-fusion, and gated-fusion representations succeed or fail under two real occlusion types. The result is deliberately conditional. Masks favor periocular specialization in AUC, sunglasses penalize it, and external transfer does not establish universal superiority of the learned gate.

This positioning also limits what can be claimed from MFR2. Because the audit found one ROF training identity also present in MFR2 (18 of 820 pairs), its results are reported with a contamination caveat and are not a fully identity-disjoint comparison with prior work. They remain useful as an implementation and domain-shift stress test, while the person-disjoint ROF protocols carry the primary evidential weight.

\subsection{Failure modes and practical implications}
The ablation exposes three failure modes that are easy to hide in a single aggregate score. The first is \emph{region failure}. Periocular LoRA is effective for masked probes because the mask leaves most of its input region visible, but its sunglasses AUC falls well below the whole-face alternatives. This is not a contradiction: specialization is useful only while the specialized region remains informative. A deployment that routes every occluded image to a periocular model would therefore be brittle to upper-face occlusion.

The second is \emph{fusion mismatch}. Fixed fusion is intentionally conservative: it always combines both normalized branches with equal weight. Learned gating has more capacity and can exploit interactions between the two embeddings, but that capacity also creates another distribution-dependent component. Its strong sunglasses and neutral AUCs show that adaptive fusion can be useful, while its exploratory RMFRD result shows that adaptation learned on one data distribution need not transfer as the best choice to another. The appropriate interpretation is therefore not that fixed fusion is intrinsically safer or that gating is intrinsically better, but that adaptive weighting must be validated on the operating conditions for which it will be used.

The third is \emph{metric failure through protocol resolution}. Low-FAR reporting is important for biometric verification, yet a target FAR can be finer than the empirical resolution of a small impostor set. In that case, the numerical TAR still follows from the implemented threshold rule, but tiny changes in one or two impostor scores can alter the estimate substantially. This is why the masked FAR-$10^{-3}$ value is retained for transparency but explicitly excluded from strong comparative claims. The much larger neutral and RMFRD impostor sets support finer empirical FAR increments, although RMFRD has separate data-integrity limitations.

These observations suggest a deployment-oriented decision process. A system should first identify the relevant occlusion regime and security operating point, then validate the candidate representation on a sufficiently large, identity-independent population from that regime. The present experiments support the use of parameter-efficient adaptation when training data and compute are limited, but they do not justify transferring a single threshold or fusion policy across datasets without calibration. They also motivate a future gate analysis that records the gate output jointly with occlusion type, genuine/impostor status, branch scores, and crop quality. Such an analysis would test whether the learned scalar behaves consistently with branch usefulness rather than merely improving the downstream loss.

\section{Limitations and Threats to Validity}
The masked ROF protocol has only nine eligible test identities and 58 genuine comparisons. Its low-FAR estimates are therefore discrete and have limited precision. MFR2 contains 820 complete-case pairs and is more balanced, but remains small for strong claims at FAR $10^{-3}$. The reproducibility archive contains multiple protocol versions; the files named in the Result Provenance section are designated as authoritative and are frozen for submission.

RMFRD has a more serious limitation. The audit reports exact duplicate image paths under different nominal identities and incomplete quality descriptors. Such records can create invalid impostor comparisons. We therefore label RMFRD exploratory and do not use it to claim clean cross-dataset generalization. A stronger revision should remove duplicate and perceptual-duplicate records, rerun all five models on the same cleaned pair set, and report the exclusion list.

The main ROF and neutral results are reported over three training seeds (7, 42, 123) with mean and standard deviation, and identity-aware paired bootstrap intervals are used for model differences. Three seeds reduce dependence on a single initialization but do not characterize every source of training variance. The ROF masked protocol is particularly small, so its low-FAR behavior should be read as descriptive rather than as a precise deployment estimate. MFR2 results are retained as an identity-overlap stress test and are not interpreted as clean cross-identity generalization. RMFRD results are exploratory due to the duplicate-path audit.

There are also construct-validity limits. The four auxiliary gate inputs are not uniformly informative across datasets; some are constant or unavailable, especially in external evaluation. The learned gate can still condition on the branch embeddings, but the present experiments do not demonstrate calibrated quality estimation. Likewise, the periocular crop is one operational definition of the eye-region context rather than a proof that the selected crop is optimal. A crop-margin ablation and an embedding-only gate would separate these effects more cleanly.

Finally, the study does not establish demographic fairness, presentation-attack resistance, or operational deployment safety. No demographic fairness conclusion is possible without subgroup labels and a prespecified subgroup protocol. The system verifies identity similarity and is not a presentation-attack detector. Thresholds appropriate for a real deployment would require a substantially larger, application-matched impostor population and an explicit cost model.
\section{Reproducibility and Artifact Specification}
The accompanying artifact contains code, raw score CSV files, metric JSON files, manifest files, generated figures, and recorded hashes. The published bundle ships the \emph{hashes} that bind every result to its checkpoint and protocol pair list, but not the checkpoint binaries themselves, the , or the face images, which are not redistributable. Consequently the artifact supports independent recomputation of result aggregation, table values, protocol construction, and provenance verification, but not a bit-exact regeneration of the trained embeddings. For complete archival reproducibility, the package retains SHA-256 values for every checkpoint and manifest, together with the exact environment lock file, dataset URLs and licenses, crop-failure counts, and the run configuration: split seed, training seed, LoRA rank, optimizer, learning rate, epoch count, hardware, and runtime. Downloaded face images should not be redistributed unless the dataset license explicitly permits redistribution.

The exact numerical provenance for the tables is: the frozen ROF pair lists and seed-sweep metric JSONs/raw-score CSVs for masked, sunglasses, and neutral evaluation; the corrected image-keyed MFR2 seed-sweep outputs; and the RMFRD-v2 metric JSONs/raw-score CSVs for exploratory transfer. Smoke-test checkpoints and partial runs are excluded from the paper tables. Metric metadata distinguish the training seed from the bootstrap RNG seed, and checkpoint/pair-list SHA-256 hashes bind each result to the evaluated model and protocol.

For independent reproduction, the intended sequence is: (1) build the normalized person-disjoint manifest; (2) generate periocular crops and record failures; (3) train M0--M4 separately for each of seeds 7, 42, and 123; (4) evaluate every checkpoint on the same frozen pair lists; (5) save raw scores before computing summary metrics; and (6) regenerate tables and figures from those saved scores. This ordering prevents figure-generation code or manuscript edits from becoming an untracked source of numerical results.
\section{Conclusion}
We studied parameter-efficient periocular fusion for identity-disjoint face verification under masks and sunglasses. LoRA adaptation and periocular/fusion representations improve selected verification measures relative to a frozen whole-face baseline. No single variant dominates. Periocular LoRA gives the best ROF masked mean AUC (0.8989$\pm$0.008), and learned gating gives the best ROF sunglasses mean AUC (0.8998$\pm$0.005) and the best neutral mean AUC (0.9441$\pm$0.002). Fixed fusion gives the highest recorded ROF masked TAR at FAR $10^{-3}$ (0.2069$\pm$0.130), though that operating point lies below the protocol FAR resolution, and fixed fusion is strongest by mean AUC in the exploratory RMFRD v2 results (0.7184$\pm$0.011). Main ROF and neutral results are reported over three training seeds (7, 42, 123) with mean$\pm$std, and identity-aware paired bootstrap intervals are provided for model differences. On MFR2, one ROF training identity is present in the test set (18 of 820 pairs, 2.2\%), so multi-seed results (mean AUC 0.8405--0.9554) are reported with that contamination caveat rather than as fully unseen-identity generalization; RMFRD results are exploratory due to the duplicate-path audit. These results support a conditional claim about adaptive fusion rather than universal robustness. The next decisive steps are direct gate-behavior analysis and a cleaned RMFRD protocol with duplicate and perceptual-duplicate exclusions.

\section{Acknowledgments}
The authors used OpenAI ChatGPT (GPT) as a writing and editing assistant for language refinement, clarity, and minor rephrasing in specific sections (abstract, introduction phrasing, and figure captions). All technical content, experiments, analysis, code, and conclusions were conceived and verified by the authors. No generative AI tool generated experimental results, figures, or references. The authors take full responsibility for the content.

\bibliographystyle{IEEEtran}
\begin{thebibliography}{00}
\bibitem{dinov2} M. Oquab \textit{et al.}, ``DINOv2: Learning robust visual features without supervision,'' \textit{Transactions on Machine Learning Research}, 2024.
\bibitem{dino} M. Caron, H. Touvron, I. Misra, H. J\'egou, J. Mairal, P. Bojanowski, and A. Joulin, ``Emerging properties in self-supervised vision transformers,'' in \textit{Proc. ICCV}, 2021, pp. 9650--9660.
\bibitem{mae} K. He, X. Chen, S. Xie, Y. Li, P. Doll\'ar, and R. Girshick, ``Masked autoencoders are scalable vision learners,'' in \textit{Proc. CVPR}, 2022, pp. 16000--16009.
\bibitem{simclr} T. Chen, S. Kornblith, M. Norouzi, and G. Hinton, ``A simple framework for contrastive learning of visual representations,'' in \textit{Proc. ICML}, 2020, pp. 1597--1607.
\bibitem{clip} A. Radford, J. W. Kim, C. Hallacy, A. Ramesh, G. Goh, S. Agarwal, G. Sastry, A. Askell, P. Mishkin, J. Clark, G. Krueger, and I. Sutskever, ``Learning transferable visual models from natural language supervision,'' in \textit{Proc. ICML}, 2021, pp. 8748--8763.
\bibitem{lora} E. J. Hu, Y. Shen, P. Wallis, Z. Allen-Zhu, Y. Li, S. Wang, L. Wang, and W. Chen, ``LoRA: Low-rank adaptation of large language models,'' in \textit{Proc. ICLR}, 2022.
\bibitem{vpt} M. Jia, L. Tang, B.-C. Chen, C. Cardie, S. Belongie, B. Hariharan, and S.-N. Lim, ``Visual prompt tuning,'' in \textit{Proc. ECCV}, 2022, pp. 709--727.
\bibitem{prompt} B. Lester, R. Al-Rfou, and N. Constant, ``The power of scale for parameter-efficient prompt tuning,'' in \textit{Proc. EMNLP}, 2021, pp. 3045--3059.
\bibitem{sphereface} W. Liu, Y. Wen, Z. Yu, M. Li, B. Raj, and L. Song, ``SphereFace: Deep hypersphere embedding for face recognition,'' in \textit{Proc. IEEE CVPR}, 2017, pp. 212--220.
\bibitem{cosface} H. Wang, Y. Wang, Z. Zhou, X. Ji, D. Gong, J. Zhou, Z. Li, and W. Liu, ``CosFace: Large margin cosine loss for deep face recognition,'' in \textit{Proc. IEEE CVPR}, 2018, pp. 5265--5274.
\bibitem{arcface} J. Deng, J. Guo, N. Xue, and S. Zafeiriou, ``ArcFace: Additive angular margin loss for deep face recognition,'' in \textit{Proc. IEEE CVPR}, 2019, pp. 4690--4699.
\bibitem{magface} Q. Meng, S. Zhao, Z. Huang, and F. Zhou, ``MagFace: A universal representation for face recognition and quality assessment,'' in \textit{Proc. IEEE CVPR}, 2021, pp. 14225--14234.
\bibitem{adaface} M. Kim, A. K. Jain, and X. Liu, ``AdaFace: Quality adaptive margin for face recognition,'' in \textit{Proc. IEEE CVPR}, 2022, pp. 18750--18759.
\bibitem{rof} M. E. Erakın, U. Demir, and H. K. Ekenel, ``On recognizing occluded faces in the wild,'' in \textit{Proc. Int. Conf. Biometrics Special Interest Group (BIOSIG)}, 2021, doi: 10.1109/BIOSIG52210.2021.9548293.
\bibitem{mfr2} A. Anwar and A. Raychowdhury, ``Masked face recognition for secure authentication,'' arXiv:2008.11104, 2020. [Online]. Available: \url{https://sites.google.com/view/masktheface/mfr2-dataset}
\bibitem{rmfrd} Z. Wang, B. Huang, G. Wang, P. Yi, and K. Jiang, ``Masked face recognition dataset and application,'' \textit{IEEE Trans. Biometrics, Behavior, and Identity Science}, vol. 5, no. 2, pp. 298--304, 2023, doi: 10.1109/TBIOM.2023.3242085.
\bibitem{llecnn} S. Ge, J. Li, Q. Ye, and Z. Luo, ``Detecting masked faces in the wild with LLE-CNNs,'' in \textit{Proc. IEEE CVPR}, 2017, pp. 2682--2690.
\bibitem{from2022} H. Qiu, D. Gong, Z. Li, W. Liu, and D. Tao, ``End2End occluded face recognition by masking corrupted features,'' \textit{IEEE Trans. Pattern Anal. Mach. Intell.}, vol. 44, no. 10, pp. 6939--6952, 2022, doi:10.1109/TPAMI.2021.3098962.
\bibitem{masklearn} L. Song, D. Gong, Z. Li, C. Liu, and W. Liu, ``Occlusion robust face recognition based on mask learning with pairwise differential Siamese network,'' in \textit{Proc. IEEE/CVF ICCV}, 2019, pp. 773--782.
\bibitem{umr} F. Boutros, N. Damer, F. Kirchbuchner, and A. Kuijper, ``Unmasking face embeddings by self-restrained triplet loss for accurate masked face recognition,'' \textit{Pattern Recognition}, vol. 124, Art. 108473, 2022.
\bibitem{focusface} P. C. Neto, F. Boutros, J. R. Pinto, N. Damer, A. F. Sequeira, and J. S. Cardoso, ``FocusFace: Multi-task contrastive learning for masked face recognition,'' in \textit{Proc. IEEE Int. Conf. Automatic Face and Gesture Recognition (FG)}, 2021, pp. 1--8.
\bibitem{beyondmasks} P. C. Neto, J. R. Pinto, F. Boutros, N. Damer, A. F. Sequeira, and J. S. Cardoso, ``Beyond masks: On the generalization of masked face recognition models to occluded face recognition,'' \textit{IEEE Access}, vol. 10, pp. 86222--86233, 2022, doi: 10.1109/ACCESS.2022.3199014.
\bibitem{hsu2022} G.-S. J. Hsu, H.-Y. Wu, C.-H. Tsai, S. Yanushkevich, and M. L. Gavrilova, ``Masked face recognition from synthesis to reality,'' \textit{IEEE Access}, vol. 10, pp. 37938--37952, 2022, doi: 10.1109/ACCESS.2022.3160828.
\bibitem{mfr2021} F. Boutros et al., ``MFR 2021: Masked face recognition competition,'' in \textit{Proc. IEEE Int. Joint Conf. Biometrics (IJCB)}, 2021, pp. 1--10, doi: 10.1109/IJCB52358.2021.9484337.
\bibitem{insightfacemfr} J. Deng, J. Guo, X. An, Z. Zhu, and S. Zafeiriou, ``Masked face recognition challenge: The InsightFace track report,'' in \textit{Proc. IEEE/CVF ICCVW}, 2021, pp. 1437--1444.
\bibitem{surveyocc} D. Zeng, R. Veldhuis, and L. Spreeuwers, ``A survey of face recognition techniques under occlusion,'' \textit{IET Biometrics}, vol. 10, no. 6, pp. 581--606, 2021.
\bibitem{periocular} U. Park, R. R. Jillela, A. Ross, and A. K. Jain, ``Periocular biometrics in the visible spectrum,'' \textit{IEEE Trans. Inf. Forensics Security}, vol. 6, no. 1, pp. 96--106, 2011, doi: 10.1109/TIFS.2010.2096810.
\bibitem{fusionRossJain} A. Ross and A. K. Jain, ``Information fusion in biometrics,'' \textit{Pattern Recognition Letters}, vol. 24, no. 13, pp. 2115--2125, 2003.
\bibitem{quality} F. Alonso-Fernandez, J. Fierrez, and J. Ortega-Garcia, ``Quality measures in biometric systems,'' \textit{IEEE Security \& Privacy}, vol. 10, no. 6, pp. 52--62, 2012, doi: 10.1109/MSP.2011.178.
\bibitem{adaptformer} S. Chen, C. Ge, Z. Tong, J. Wang, Y. Song, J. Wang, and P. Luo, ``AdaptFormer: Adapting vision transformers for scalable visual recognition,'' in \textit{Advances in Neural Information Processing Systems}, vol. 35, 2022.
\bibitem{fbr2024} L. C. O. Tiong, D. Sigmund, C.-H. Chan, and A. B. J. Teoh, ``Flexible biometrics recognition: Bridging the multimodality gap through attention alignment and prompt tuning,'' in \textit{Proc. IEEE/CVF Conf. Computer Vision and Pattern Recognition (CVPR)}, 2024, pp. 267--276.
\bibitem{ckd2024} Y. G. Jung, J. Park, C. Y. Low, J. C. L. Chai, L. C. O. Tiong, and A. B. J. Teoh, ``Periocular embedding learning with consistent knowledge distillation from face,'' \textit{Neurocomputing}, vol. 575, Art. 127263, 2024, doi: 10.1016/j.neucom.2024.127263.
\end{thebibliography}
\end{document}
